\section{Results}
We present results from running the adversarial scenarios in our approach for each of the research questions in the Sections below. 

\subsection{RQ1: Impact of Adversarial Generation on Threat Severity}
\begin{table*}[t]
\centering
\caption{Mean threat metric values across 30 scenarios per typology for seed (baseline near-miss) and generated (adversarial) scenarios using Baidu Apollo. Lower TTC, THW, and LTTC indicate greater risk; higher BTN and STN reflect greater risk. Bold values breach the emergency threshold (Table~\ref{tab:thresholds}).}
\label{tab:rq1_threat_comparison}
\renewcommand{\arraystretch}{1.0}
\begin{tabular}{@{}l cc cc cc cc cc@{}}
\toprule
& \multicolumn{2}{c}{\textbf{TTC [s] $\downarrow$}} & \multicolumn{2}{c}{\textbf{THW [s] $\downarrow$}} & \multicolumn{2}{c}{\textbf{BTN [--] $\uparrow$}} & \multicolumn{2}{c}{\textbf{LTTC [s] $\downarrow$}} & \multicolumn{2}{c}{\textbf{STN [--] $\uparrow$}} \\
\cmidrule(lr){2-3} \cmidrule(lr){4-5} \cmidrule(lr){6-7} \cmidrule(lr){8-9} \cmidrule(lr){10-11}
\textbf{Typology} & Seed & Gen. & Seed & Gen. & Seed & Gen. & Seed & Gen. & Seed & Gen. \\
\midrule
\rowcolor{blue!5}
(1) Lead Veh. Stopped         & 1.60 & \textbf{1.35} & 0.80 & \textbf{0.68} & 0.781 & \textbf{0.82} & -- & -- & -- & -- \\
\rowcolor{blue!5}
(2) Lead Veh. Decelerating    & 1.80 & \textbf{1.54} & 0.90 & \textbf{0.72} & 0.617 & 0.63 & -- & -- & -- & -- \\
\rowcolor{blue!5}
(3) Lead Veh. Lower Speed     & 2.40 & \textbf{1.54} & 0.60 & 0.77* & 0.174 & 0.63 & -- & -- & -- & -- \\
\rowcolor{green!5}
(4) Turning at Non-Signalized  & 1.54 & \textbf{1.20} & -- & -- & 0.63 & 0.60 & -- & -- & -- & -- \\
\rowcolor{green!5}
(5) Straight Crossing          & 1.80 & \textbf{1.20} & -- & -- & 0.46 & 0.70 & -- & -- & -- & -- \\
\rowcolor{green!5}
(6) LTAP/OD Signalized         & 2.10 & \textbf{1.44} & -- & -- & 0.322 & 0.68 & -- & -- & -- & -- \\
\rowcolor{green!5}
(7) LTAP/OD Non-Signalized    & 1.92 & \textbf{1.20} & -- & -- & 0.43 & 0.69 & -- & -- & -- & -- \\
\rowcolor{orange!5}
(8) Changing Lanes             & -- & -- & -- & -- & -- & -- & 1.40 & \textbf{0.48} & 0.279 & 0.689 \\
\rowcolor{orange!5}
(9) Turning Same Dir.          & -- & -- & -- & -- & -- & -- & 1.50 & \textbf{0.50} & 0.126 & 0.625 \\
\rowcolor{orange!5}
(10) Drifting Same Dir.        & -- & -- & -- & -- & -- & -- & 0.73 & \textbf{0.38} & 0.320 & \textbf{1.056} \\
\bottomrule
\end{tabular}
\end{table*}

Table~\ref{tab:rq1_threat_comparison} compares mean threat metric values between seed and generated scenarios for the different adversarial typologies. We show this only for Apollo due to space constraints; the corresponding table for Traffic Manager is available in the anonymous repository.  We find that our approach is successful in consistently increasing threat severity across all ten adversarial typologies, with one exception noted below. Note that TTC, THW, and LTTC are inversely related to risk with lower values indicating a more safety-critical situation.  Whereas BTN and STN are directly proportional to risk, with higher values reflecting greater threat severity. 

\paragraph{Longitudinal Conflicts.}
TTC decreases across all three car-following typologies, most substantially in typology~3 (Lead Vehicle at Lower Speed): from 2.40\,s to 1.54\,s (36\% reduction), falling below the 1.5\,s emergency threshold. THW decreases for typologies~1 and~2, both below the 1.0\,s emergency threshold. The sole exception is typology~3, where THW increases slightly from 0.60\,s to 0.77\,s, as the EA prioritises TTC and BTN reduction when THW is already well below emergency threshold. BTN increases across all three typologies; for typology~1 it reaches 0.82, exceeding the 0.80 emergency threshold.

\paragraph{Junction Conflicts.}
All four intersection typologies show TTC reductions below 1.44\,s, with the largest in typologies~5 and~7 (both reaching 1.20\,s). BTN increases substantially, with typology~6 (LTAP/OD Signalized) showing the largest absolute increase from 0.32 to 0.68, indicating that the EA effectively tightens intersection timing to create severe braking demands.

\paragraph{Lateral Conflicts.}
The lateral typologies exhibit the most pronounced increases. LTTC drops by 48--67\% across all three typologies, with all generated values at or below the 0.5\,s emergency threshold. STN increases correspondingly; typology~10 (Drifting) reaches 1.056, exceeding 1.0 and indicating that the required lateral evasion surpasses the vehicle's physical capacity.

\paragraph{Summary.}
Our approach successfully increases threat severity across all ten adversarial typologies. Lateral conflicts are most affected, with LTTC reductions exceeding 50\% and STN approaching the vehicle's physical evasion capacity. Junction conflicts show the largest absolute BTN increases, and longitudinal conflicts exhibit consistent but more moderate degradation, with one minor exception: THW increases slightly for typology~3, as the algorithm prioritises other metrics already well below the emergency threshold.

\subsection{RQ2: Collision and Failure Outcomes}
\begin{table*}[t]
    \caption{Failure counts recorded for CARLA Traffic Manager (TM) and Baidu Apollo across kinematic scenario categories in the generated adversarial scenarios. Each cell reports the total number of failures (collisions and failure to reach destination), with the number of destination failures shown in parentheses.}
  \label{tab:collision_comparison_all}
  \renewcommand{\arraystretch}{1.0} 
  \begin{tabularx}{\textwidth}{@{} 
      >{\hsize=0.8\hsize\raggedright\arraybackslash}X 
      >{\hsize=1.6\hsize\raggedright\arraybackslash}X 
      >{\hsize=0.8\hsize\centering\arraybackslash}X 
      >{\hsize=0.8\hsize\centering\arraybackslash}X 
      @{}}
    \toprule
    \textbf{Kinematic Category} & \textbf{Typology} & \multicolumn{2}{c}{\textbf{\# collisions and undesirable behaviour}} \\
    \cmidrule(l){3-4}
    & & \textbf{TM} & \textbf{Apollo} \\
    \midrule
    \rowcolor{blue!5}
    \textbf{Longitudinal} & 
    (1) Lead Vehicle Stopped & 4(\textcolor{blue}{1}) $\uparrow$ & 1(\textcolor{blue}{1}) $\uparrow$ \\
    \rowcolor{blue!5}
     & 
    (2) Lead Vehicle Decelerating & 1 $\uparrow$ & 0 \\
    \rowcolor{blue!5}
    & (3) Lead Veh. Moving at Lower Speed & 0 & 0 \\
    \rowcolor{green!5}
    \textbf{Junction} & 
    (4) Veh. Turning at Non-Signalized & 24(\textcolor{blue}{6}) $\uparrow$ & 3(\textcolor{blue}{1}) $\uparrow$ \\
    \rowcolor{green!5}
     & 
    (5) Straight Crossing Non-Signalized & 12(\textcolor{blue}{3}) $\uparrow$ & 3 $\uparrow$ \\
    \rowcolor{green!5}
    & (6) LTAP/OD at Signalized & 17(\textcolor{blue}{2}) $\uparrow$ & 2(\textcolor{blue}{2}) $\uparrow$ \\
    \rowcolor{green!5}
    & (7) LTAP/OD at Non-Signalized & 27(\textcolor{blue}{7}) $\uparrow$ & 6(\textcolor{blue}{2}) $\uparrow$ \\
    \rowcolor{orange!5}
    \textbf{Lateral} & 
    (8) Veh. Changing Lanes (Same Dir.) & 19(\textcolor{blue}{3}) $\uparrow$ & 5(\textcolor{blue}{1}) $\uparrow$ \\
    \rowcolor{orange!5}
     & 
    (9) Veh. Turning (Same Dir.) & 14(\textcolor{blue}{7}) $\uparrow$ & 3(\textcolor{blue}{2}) $\uparrow$ \\
    \rowcolor{orange!5}
    & (10) Vehicle Drifting (Same Dir.) & 30(\textcolor{blue}{8}) $\uparrow$ & 6(\textcolor{blue}{4}) $\uparrow$ \\
    \bottomrule
  \end{tabularx}
\end{table*}

Table~\ref{tab:collision_comparison_all} presents failure counts for both ADS on the 300 adversarial scenarios. In the original SCTrans baseline, collisions occur in only 3.3\% of scenarios. The adversarial scenarios increase failure rates to 9.7\% for Apollo (29 failures) and 49.3\% for Traffic Manager (148 failures) out of total number of 300 generated scenarios.

\paragraph{Longitudinal Conflicts.}
Despite the threat increases reported in RQ1, longitudinal scenarios produce few failures: Apollo~1, TM~5 across 90 scenarios. Even typology~3, which showed the largest relative threat increase, produces zero failures for both systems. The one-dimensional nature of car-following conflicts provides a clear braking response that absorbs elevated threat levels.

\paragraph{Junction Conflicts.}
Junction scenarios show stronger correspondence between threat increases and failures: Apollo records 14 failures across 120 scenarios, TM records 80. Typology~7 (LTAP/OD Non-Signalized) produces the highest counts (TM:~27, Apollo:~6). Both systems are particularly vulnerable at non-signalized junctions, where no traffic signals regulate conflict timing and the ADS must independently predict the POV's path. Destination failures also concentrate here, suggesting that junction perturbations disrupt route planning beyond immediate collision avoidance.

\paragraph{Lateral Conflicts.}
Lateral scenarios produce the highest per-typology failure rates. Typology~10 (Drifting) triggers failures in all 30 TM scenarios and 6 Apollo scenarios. Lateral scenarios also yield the highest proportion of destination failures --- typology~10 (8/30 TM) and typology~9 (7/14 TM) --- indicating disruption to downstream path planning.

\paragraph{Summary.}
Adversarial scenarios substantially increase failure rates compared to the 3.3\% baseline, reaching 9.7\% for Apollo and 49.3\% for TM. Lateral conflicts are the most hazardous --- typology~10 (Drifting) triggers failures in all 30 TM scenarios --- followed by junction conflicts, with non-signalised junctions (typologies~4 and~7) concentrating the most failures. Longitudinal conflicts produce the most modest increase. Destination failures are disproportionately concentrated in lateral and junction scenarios, indicating that these behaviours disrupt downstream path planning beyond immediate collision avoidance.

\subsection{RQ3: ADS Robustness Comparison}

We now compare the two ADS using the results from Tables~\ref{tab:rq1_threat_comparison} and~\ref{tab:collision_comparison_all}. Overall, Traffic Manager's failure rate (49.3\%) is five times Apollo's (9.7\%), with the gap varying by kinematic category.

In \textit{longitudinal conflicts}, the difference is minimal (TM:~5, Apollo:~1), as both architectures handle one-dimensional braking effectively. In \textit{junction conflicts}, the gap widens to six-fold (TM:~80 vs Apollo:~14). Apollo's learned prediction modules provide a clear advantage in multi-directional conflicts requiring trajectory reasoning through intersections, whereas TM's fixed priority rules cannot adapt to tightened adversarial timing. Destination failures are disproportionately concentrated in TM's junction results (18 of 80). In \textit{lateral conflicts}, TM fails in 63 of 90 scenarios (70.0\%) versus Apollo's 14 (15.6\%). Typology~10 is catastrophic for TM (30/30) but Apollo is robust in the majority of cases, demonstrating that its perception pipeline can detect gradual encroachment. Typology~9 shows a distinctive pattern where 50\% of TM's failures are destination failures, indicating that lateral disruption derails path-planning even when immediate collision is avoided.

\paragraph{Summary.}
TM's rule-based architecture is consistently more vulnerable than Apollo's ML-based pipeline, with the robustness gap scaling with the dimensionality of the conflict: minimal for longitudinal braking, six-fold for multi-directional junction reasoning, and 4.5-fold for lateral evasion. Destination failures concentrate almost exclusively in TM, particularly under lateral and junction perturbations, suggesting that rule-based planning degrades not only in immediate collision avoidance but also in longer-horizon route recovery. Notably, both architectures remain vulnerable to gradual lateral encroachment and non-signalised junction conflicts, indicating that these scenario types warrant safety improvements regardless of the underlying ADS design.

